\section{Full ELBO derivation}
\label{app:elbo}

This appendix gives the full derivation of the ELBO summarized in Section~\ref{sec:elbo}. We start from the marginal log-likelihood, apply Jensen's inequality, substitute the joint factorization and variational family, decompose into reconstruction and KL terms, write the closed-form Gaussian-Gaussian KLs out in full, and treat the flow KL by Monte Carlo with the change-of-variables formula.

\subsection*{Starting point}
Let $\bm Z=(\bm\mu^{\mathrm{pop}},\bdelta,\bm\eta)$ collect all latents on the logit scale, with $\bm\mu^{\mathrm{pop}}=(\mu^{\mathrm{pop}}_1,\dots,\mu^{\mathrm{pop}}_L)$, $\bdelta=(\delta_1,\dots,\delta_S)$, $\bm\eta=(\eta_{s,\ell})_{s,\ell}$. The marginal log-likelihood of the observed fragments $F=(F_{s,\ell})_{s,\ell}$ is
\begin{equation}
\log p(F) = \log \int p(F,\bm Z)\,d\bm Z.
\end{equation}
Multiplying and dividing by $q(\bm Z)$ and applying Jensen's inequality to the concave $\log$ gives
\begin{equation}
\log p(F) = \log \E_q\!\left[\frac{p(F,\bm Z)}{q(\bm Z)}\right] \ge \E_q\!\left[\log\frac{p(F,\bm Z)}{q(\bm Z)}\right] \;=\; \elbo(q),
\label{eq:appA-elbo-def}
\end{equation}
with equality if and only if $q(\bm Z)=p(\bm Z\mid F)$ almost everywhere, since the gap is exactly $\KL(q\|p(\cdot\mid F))$ (Proposition~\ref{prop:elbo}).

\subsection*{Substituting the joint factorization}
The joint $p(F,\bm Z)$ factorizes as in \eqref{eq:joint}:
\begin{equation}
p(F,\bm Z) = \prod_\ell p(\mu^{\mathrm{pop}}_\ell)\,\prod_s p(\delta_s)\,\prod_{s,\ell} p(\eta_{s,\ell}\mid \mu^{\mathrm{pop}}_\ell,\delta_s)\,p(F_{s,\ell}\mid\eta_{s,\ell}),
\end{equation}
and the variational family $q$ factorizes as in \eqref{eq:varfamily}:
\begin{equation}
q(\bm Z) = \prod_\ell q_{\lambda_\ell}(\mu^{\mathrm{pop}}_\ell)\,\prod_s q_{\nu_s}(\delta_s)\,\prod_{s,\ell} q_\phi(\eta_{s,\ell}\mid F_{s,\ell};m_\ell,m_s^\delta).
\end{equation}
Substituting both into \eqref{eq:appA-elbo-def} and grouping log-ratios by latent gives
\begin{align}
\elbo(q) =\;& \sum_{s,\ell}\E_q[\log p(F_{s,\ell}\mid \eta_{s,\ell})]
\;-\;\sum_\ell \KL\big(q_{\lambda_\ell}(\mu^{\mathrm{pop}}_\ell)\,\big\|\,p(\mu^{\mathrm{pop}}_\ell)\big) \nonumber\\
&-\;\sum_s \KL\big(q_{\nu_s}(\delta_s)\,\big\|\,p(\delta_s)\big)
\;-\;\sum_{s,\ell}\E_{q_{\lambda_\ell} q_{\nu_s}}\!\big[\KL\big(q_\phi(\eta_{s,\ell}\mid \cdot)\,\big\|\,p(\eta_{s,\ell}\mid\mu^{\mathrm{pop}}_\ell,\delta_s)\big)\big].
\label{eq:appA-decomp}
\end{align}
This is the decomposition stated in \eqref{eq:elbo-decomp}.

\subsection*{Reconstruction terms}
The fragment likelihood factorizes over three observation modalities:
\begin{equation}
\log p(F_{s,\ell}\mid\eta_{s,\ell}) = \log\BB(n^m_{s,\ell};n^{\mathrm{cpg}}_{s,\ell},\kappa\sigm(\eta),\kappa(1-\sigm(\eta))) + \sum_j \log\Cat(m_{s,\ell,j};\bpi(\sigm(\eta),c^{\mathrm{seq}}_\ell)) + \log\NB(n_{s,\ell};r_\ell,p(\sigm(\eta),\mathrm{GC}_\ell)).
\end{equation}
Each is a function of $\eta_{s,\ell}$, and the expectation $\E_{q_\phi}[\log p(F\mid\eta)]$ is computed by Monte Carlo over flow samples $\eta=T_\phi(\epsilon;c_{s,\ell})$.

\subsection*{Closed-form Gaussian-Gaussian KLs}
For two univariate Gaussians,
\begin{equation}
\KL(\N(m,v)\,\|\,\N(\mu,\tau^2)) = \tfrac{1}{2}\!\left[\frac{v}{\tau^2} + \frac{(m-\mu)^2}{\tau^2} - 1 - \log\frac{v}{\tau^2}\right].
\label{eq:appA-gauss-kl}
\end{equation}
Applied to the population layer with $q_{\lambda_\ell}=\N(m_\ell,v_\ell)$ and $p=\N(\mu_0,\tau_0^2)$, this gives the locus-wise KL contribution. Applied to the sample-shift layer with $q_{\nu_s}=\N(m_s^\delta,v_s^\delta)$ and $p=\N(0,\sigma_\delta^2)$, the analogous formula obtains.

\subsection*{Flow KL by Monte Carlo}
The conditional KL $\KL(q_\phi(\eta\mid c)\|p(\eta\mid\mu^{\mathrm{pop}},\delta))$ has no closed form because $q_\phi$ is a normalizing flow. We estimate it by single-sample Monte Carlo:
\begin{equation}
\widehat{\KL} = \log q_\phi(\eta\mid c) - \log p(\eta\mid \mu^{\mathrm{pop}},\delta),\quad \eta = T_\phi(\epsilon;c),\;\epsilon\sim\N(0,1),
\end{equation}
where the flow log-density is evaluated by the change-of-variables formula \eqref{eq:flow-density}:
\begin{equation}
\log q_\phi(\eta\mid c) = \log\N(\epsilon;0,1) - \log|\det\nabla_\epsilon T_\phi(\epsilon;c)|.
\end{equation}

\subsection*{Tightness}
Equality in \eqref{eq:appA-elbo-def} holds iff $q=p(\cdot\mid F)$. Within the structured family of Section~\ref{sec:varfamily}, this is achievable only if the flow $T_\phi$ is sufficiently expressive to capture the conditional posterior, the population layer is exactly Gaussian (it is, by construction in conjugate cases), and the mean-field assumption across plates is exactly satisfied. In practice, the gap between the family and the true posterior is reflected in the amortization gap (Section~\ref{ssec:amortization}) and in the IWAE-tightening improvements of Section~\ref{ssec:iwae}.
